%% Boltzmannator — theory and user guide
\documentclass[11pt,a4paper]{article}

\usepackage{amsmath,amssymb,mathtools}
\usepackage[margin=2.5cm]{geometry}
\usepackage{hyperref}
\usepackage{xcolor}
\usepackage{booktabs}
\usepackage{parskip}
\usepackage{microtype}
\usepackage{authblk}

\hypersetup{colorlinks=true, linkcolor=blue!60!black, urlcolor=blue!60!black}

\newcommand{\E}{\mathbb{E}}
\newcommand{\R}{\mathbb{R}}
\newcommand{\pz}{p_z}
\newcommand{\px}{p_x}
\newcommand{\ft}{f_{\!\theta}}
\newcommand{\Jf}{J_{\!f}}
\newcommand{\kT}{k_{\!B}T}
\newcommand{\param}{\theta}
\newcommand{\sigmoid}{\sigma}

\title{\textbf{Boltzmannator}\\[4pt]
       \large {\em A Visual Explorer for One-Dimensional Normalising Flows}}



\author[1]{Claude Code}

\author[2]{Christoph Dellago}

\affil[1]{Anthropic}

\affil[2]{Faculty of Physics, University of Vienna}

\date{}

\begin{document}
\maketitle
\tableofcontents
\bigskip

%% ─────────────────────────────────────────────────────────────────────────────
\section{Introduction}

\textbf{Boltzmannator} is an interactive, browser-based application for
exploring normalising flows in one dimension.  Given a user-chosen
\emph{latent} distribution $\pz(z)$ and a parametric transformation
$x = \ft(z)$, the app visualises how the probability density is reshaped
under the map, computes the push-forward density $\px(x)$, and trains the
parameters $\param$ against either a target Boltzmann distribution
(energy-based mode) or a set of example data points (maximum-likelihood
mode).  The aim is to make the interplay between the transformation, the
Jacobian, and the resulting density immediately tangible.

Normalising flows were introduced as a tool for flexible probability-density
modelling and variational inference \cite{rezende2015}; see
\cite{papamakarios2021} for a comprehensive review.  In statistical physics
the same idea underlies \emph{Boltzmann Generators}, which learn a flow that
maps a simple latent density onto a target Boltzmann distribution to sample
equilibrium states of many-body systems \cite{noe2019boltzmann} (see
\cite{coretti2024kimreview} for a short review).  Boltzmann Generators have
since been extended to free-energy estimation via learned mappings
\cite{wirnsberger2020}, to rare-event and transition-path sampling
\cite{falkner2023}, to the efficient mapping of phase diagrams and equilibrium
sampling of large-scale materials with conditional flows
\cite{schebek2024,schebek2025}, and to learning mappings between equilibrium
states of liquid systems \cite{coretti2025jcp}.  Boltzmannator illustrates
these ideas in the simplest possible one-dimensional setting.

The program is written in Python using the \href{https://nicegui.io}{NiceGUI}
framework, with all numerics in NumPy and the figures rendered by Matplotlib.
It runs as a small local web server and is used through a web browser at
\texttt{http://localhost:8080}.  Because it is served over the network, several
users on the same local network can connect simultaneously (e.g.\ students in a
classroom opening \texttt{http://\textit{host-ip}:8080}); each browser gets its
own fully independent session, so one user's parameters and training never
affect another's.  Installation and start-up instructions are given in the
accompanying \texttt{README}.

%% ─────────────────────────────────────────────────────────────────────────────
\section{Mathematical Background}

\subsection{The Change-of-Variables Formula}

Let $z \in \R$ be a random variable with density $\pz(z)$, and let
$x = \ft(z)$ be a differentiable, invertible (monotone) function
parametrised by $\param$.  The \emph{push-forward} density of $x$ is
obtained from the change-of-variables formula:
\begin{equation}
  \px(x) \;=\; \frac{\pz\!\left(\ft^{-1}(x)\right)}{|\Jf(z)|}
  \;=\; \frac{\pz(z)}{|\Jf(z)|},
  \label{eq:cov}
\end{equation}
where $z = \ft^{-1}(x)$ is the pre-image and
\begin{equation}
  \Jf(z) \;=\; \frac{d\ft}{dz}
\end{equation}
is the (scalar) Jacobian.  Equation~\eqref{eq:cov} is the cornerstone of
normalising flows: by choosing $\ft$ appropriately, any latent density
$\pz$ can be mapped to an arbitrarily complex target density $\px$.

The app displays the four quantities in \eqref{eq:cov} simultaneously:
$\pz(z)$ (top strip), the curve $x = \ft(z)$ (main panel),
$|\Jf(z)|^{-1}$ (bottom strip), and $\px(x)$ (right panel).
The bottom strip shows the \emph{inverse} Jacobian because
$\px(x) = \pz(z)\times|\Jf(z)|^{-1}$: the value directly reads off the
local density-magnification factor ($>1$ compresses, $<1$ dilutes).

\subsection{Latent Distributions}

The following latent densities are available, all parametrised by a mean
$\mu$ and a scale $\sigma$:

\begin{center}
\renewcommand{\arraystretch}{1.4}
\begin{tabular}{lll}
\toprule
Name & Density $\pz(z)$ & Standard deviation \\
\midrule
Gaussian  & $\dfrac{1}{\sigma\sqrt{2\pi}}\exp\!\left(-\dfrac{(z-\mu)^2}{2\sigma^2}\right)$
           & $\sigma$ \\
Uniform   & $\dfrac{1}{2\sigma\sqrt{3}}$ for $|z-\mu| \le \sigma\sqrt{3}$, else $0$
           & $\sigma$ \\
Laplace   & $\dfrac{1}{\sigma\sqrt{2}}\exp\!\left(-\dfrac{\sqrt{2}\,|z-\mu|}{\sigma}\right)$
           & $\sigma$ \\
Bimodal   & $\dfrac{1}{2}\mathcal{N}(z;-\mu,\sigma)
             +\dfrac{1}{2}\mathcal{N}(z;\,+\mu,\sigma)$
           & $\sqrt{\mu^2+\sigma^2}$ \\
\bottomrule
\end{tabular}
\end{center}

\noindent
For the \textbf{Bimodal} distribution the slider $\mu$ sets the
\emph{peak positions} at $\pm\mu$ (not the overall mean, which is always
zero), while $\sigma$ sets the width of each individual peak.  The slider
is therefore restricted to $\mu \ge 0$; the default values are $\mu=1$,
$\sigma=0.5$.

\subsection{The Target Boltzmann Distribution}

The target density used for energy-based training is a one-dimensional
Boltzmann distribution,
\begin{equation}
  p^*(x) \;=\; \frac{1}{Z}\exp\!\left(-\frac{U(x)}{\kT}\right),
  \qquad
  Z = \int_{-\infty}^{\infty}\exp\!\left(-\frac{U(x)}{\kT}\right)dx,
\end{equation}
with the polynomial potential energy
\begin{equation}
  U(x) \;=\; u_1 x + u_2 x^2 + u_3 x^3 + u_4 x^4.
\end{equation}
The coefficients $u_1,\ldots,u_4$ and the temperature $T$ are controlled
via sliders.  A positive $u_4$ ensures that $U \to +\infty$ for
$|x|\to\infty$ and that the partition function $Z$ is finite.

\subsubsection*{Exact transformation}

When the target distribution is set up, the theoretically perfect
transformation $x^*(z)$ can be computed as the composition of CDFs:
\begin{equation}
  x^*(z) \;=\; F_{\,p^*}^{-1}\!\bigl(F_{\pz}(z)\bigr),
\end{equation}
where $F_{\pz}$ and $F_{p^*}$ are the cumulative distribution functions of
the latent and target densities, respectively.  This \emph{exact} curve is
shown as a dashed reference line in the transformation panel.  Toggling
\emph{Show exact transformation} does not change the displayed $z$- or
$x$-range (the curve is simply drawn, clipped to the current axes).

%% ─────────────────────────────────────────────────────────────────────────────
\section{Parametric Transformations}

\subsection{Polynomial Transformation}

The polynomial family is
\begin{equation}
  \ft(z) \;=\; \theta_0 + \theta_1 z + \theta_2 z^2 + \theta_3 z^3,
  \label{eq:poly}
\end{equation}
with Jacobian
\begin{equation}
  \Jf(z) \;=\; \theta_1 + 2\theta_2 z + 3\theta_3 z^2.
\end{equation}
A cubic can represent shifts, scalings, skewness, and mild non-linearity,
but it is not guaranteed to be monotone: if $\Jf$ changes sign the
transformation is not invertible and the change-of-variables formula
\eqref{eq:cov} does not apply.  The app highlights regions of
non-monotonicity in orange and displays a warning in the main panel.

\subsection{Rational-Quadratic Spline (RQS) Transformation}
\label{sec:rqs}

The RQS family \cite{durkan2019neural} defines a piecewise
rational-quadratic bijection on a finite support $[-B, B]$ with $K$
bins, and extends linearly (identity) outside the support.

\paragraph{Parameters.}
Given $K$ bins, the raw parameters are:
\begin{itemize}
  \item Boundary $B > 0$ (one slider).
  \item Unconstrained width logits $\tilde w_1,\ldots,\tilde w_K$: actual
        bin widths $\Delta x_k = 2B \cdot \mathrm{softmax}(\tilde w)_k$,
        so $\sum_k \Delta x_k = 2B$.
  \item Unconstrained height logits $\tilde h_1,\ldots,\tilde h_K$: actual
        heights $\Delta y_k = 2B \cdot \mathrm{softmax}(\tilde h)_k$.
  \item Unconstrained derivative logits $\tilde d_0,\ldots,\tilde d_K$:
        actual knot derivatives $d_k = e^{\tilde d_k} > 0$.
        Setting $\tilde d_k = 0$ gives $d_k = 1$.
\end{itemize}
Knot positions are $x_0 = y_0 = -B$ and
$x_{k+1} = x_k + \Delta x_k$, $y_{k+1} = y_k + \Delta y_k$.
The total number of parameters is $3K + 2$ (including $B$).

\paragraph{Forward pass.}
For $z$ in bin $k$ (i.e.\ $x_k \le z < x_{k+1}$), define
$\xi = (z - x_k)/\Delta x_k \in [0,1]$ and
$s_k = \Delta y_k / \Delta x_k$.  Then
\begin{equation}
  \ft(z) \;=\; y_k
  \;+\; \Delta y_k\,\frac{s_k\,\xi^2 + d_k\,\xi(1-\xi)}
                        {s_k + (d_{k+1}+d_k-2s_k)\,\xi(1-\xi)},
  \label{eq:rqs_fwd}
\end{equation}
with Jacobian
\begin{equation}
  \Jf(z) \;=\;
  \frac{s_k^2\bigl[d_{k+1}\,\xi^2 + 2s_k\,\xi(1-\xi) + d_k(1-\xi)^2\bigr]}
       {\bigl[s_k + (d_{k+1}+d_k-2s_k)\,\xi(1-\xi)\bigr]^2}.
  \label{eq:rqs_jac}
\end{equation}
Since all $\Delta x_k, \Delta y_k, d_k > 0$ the transformation is always
strictly monotone ($\Jf > 0$ everywhere in the support).

\paragraph{Inverse pass.}
Given $x_\mathrm{out}$ in bin $k$ (located by $y_k \le x_\mathrm{out} < y_{k+1}$),
the normalised position $\xi$ is the root of the quadratic
$a\xi^2 + b\xi + c = 0$ with
\begin{align}
  a &= \Delta y_k(s_k - d_k) + (x_\mathrm{out}-y_k)(d_{k+1}+d_k-2s_k),\\
  b &= \Delta y_k\,d_k - (x_\mathrm{out}-y_k)(d_{k+1}+d_k-2s_k),\\
  c &= -s_k(x_\mathrm{out}-y_k).
\end{align}
The physically meaningful root (in $[0,1]$) is selected; the original
input is then recovered as $z = x_k + \xi\,\Delta x_k$.  Because the
inverse is analytic, it requires no iterative bisection, making
example-based training fast even for large datasets.

\paragraph{Gradient.}
Analytical gradients through the softmax/exp re-parameterisations are
complex.  The app uses central finite differences instead
($\varepsilon = 10^{-4}$), which requires $2(3K{+}2)$ forward passes
per epoch — acceptable because each pass is $O(N)$.

\subsection{Single-Layer Perceptron (SLP) Transformation}

The SLP family combines a linear part with a sum of $K \le 8$ sigmoid
units:
\begin{equation}
  \ft(z) \;=\; a + b\,z + \sum_{k=1}^{K} w_k\,\sigmoid\!\left(\frac{z-c_k}{s_k}\right),
  \label{eq:slp}
\end{equation}
where $\sigmoid(t) = (1+e^{-t})^{-1}$ is the logistic sigmoid and
$(w_k, c_k, s_k)$ are the weight, centre, and scale of the $k$-th unit.
The Jacobian is
\begin{equation}
  \Jf(z) \;=\; b + \sum_{k=1}^{K} \frac{w_k}{s_k}\,
                \sigmoid\!\left(\frac{z-c_k}{s_k}\right)
                \left[1 - \sigmoid\!\left(\frac{z-c_k}{s_k}\right)\right].
\end{equation}
Because $\sigmoid(1-\sigmoid) \ge 0$, enforcing $w_k > 0$ and $b > 0$
guarantees $\Jf(z) > 0$ everywhere, making \eqref{eq:slp} a
\emph{strictly monotone} transformation.  The SLP can represent
asymmetric, multimodal, and sigmoidal reshapings that the cubic polynomial
cannot.

The inverse $z = \ft^{-1}(x)$ has no closed form for $K \ge 1$; it is
computed numerically by bisection (see Section~\ref{sec:bisection}).

%% ─────────────────────────────────────────────────────────────────────────────
\section{Training the Flow}

\subsection{Energy-Based Training}
\label{sec:energy}

Given the target Boltzmann distribution $p^*(x)$, we want the push-forward
$\px$ to approximate $p^*$.  A natural objective is the
Kullback--Leibler divergence from the model to the target:
\begin{equation}
  D_{\mathrm{KL}}(\px \| p^*) \;=\;
  \E_{\px}\!\left[\log\frac{\px(x)}{p^*(x)}\right]
  \;=\;
  \E_{z\sim\pz}\!\left[\frac{U(\ft(z))}{\kT}\right]
  - \E_{z\sim\pz}\!\left[\log|\Jf(z)|\right]
  + \text{const},
\end{equation}
where the constant $\log Z$ is independent of $\param$.  The \emph{variational
free energy} (the trainable part) is therefore
\begin{equation}
  \boxed{
  \mathcal{L}_{\text{energy}}(\param) \;=\;
  \underbrace{\E_z\!\left[\frac{U(\ft(z))}{\kT}\right]}_{\text{energy term}}
  \;-\;
  \underbrace{\E_z\!\left[\log|\Jf(z)|\right]}_{\text{entropy term}}.
  }
  \label{eq:energy_loss}
\end{equation}
The energy term penalises configurations in high-potential regions; the
entropy term rewards transformations with large Jacobian (i.e.\ spreading
mass).  Together they drive the model towards the Boltzmann target.

\subsection{Example-Based Training (Maximum Likelihood)}
\label{sec:nll}

When example data $\{x_i\}_{i=1}^N$ from an unknown distribution
$p_{\mathrm{data}}(x)$ are available, the parameters can be trained by
maximising the log-likelihood of the data under the model, or equivalently
by minimising the negative log-likelihood (NLL):
\begin{equation}
  \boxed{
  \mathcal{L}_{\text{NLL}}(\param) \;=\;
  \E_{x\sim p_{\mathrm{data}}}\!\left[-\log\pz(z) + \log|\Jf(z)|\right],
  }
  \label{eq:nll_loss}
\end{equation}
where $z = \ft^{-1}(x)$ denotes the pre-image of each data point.  The
two terms can be interpreted as
\begin{itemize}
  \item $-\E[\log\pz(z)]$: the latent-space entropy cost — the mapped
        pre-images should be likely under $\pz$;
  \item $+\E[\log|\Jf(z)|]$: the volume-change cost — a large Jacobian
        compresses the target density, so the optimiser
        balances spreading against fitting.
\end{itemize}
Training example data can be generated from the target Boltzmann
distribution via the CDF-inversion trick described in Section~2.3.

\subsection{Analytical Gradients}

Both losses are minimised by stochastic gradient descent with analytical
gradients, avoiding the $O(d)$ extra function evaluations of finite
differences.

\paragraph{Energy-based gradient.}
For a parameter $\param_i$,
\begin{equation}
  \frac{\partial\mathcal{L}_{\text{energy}}}{\partial\param_i}
  \;=\;
  \E_z\!\left[\frac{dU}{dx}\,\frac{\partial\ft}{\partial\param_i}\,\frac{1}{\kT}
              \;-\;
              \frac{1}{\Jf(z)}\frac{\partial\Jf}{\partial\param_i}\right].
\end{equation}

\paragraph{NLL gradient.}
Using the implicit-function theorem, $\partial z/\partial\param_i =
-(\partial\ft/\partial\param_i)/\Jf(z)$, one obtains
\begin{equation}
  \frac{\partial\mathcal{L}_{\text{NLL}}}{\partial\param_i}
  \;=\;
  \E_x\!\left[
    \left(s(z) - \frac{\partial_z\Jf}{\Jf}\right)
    \frac{\partial\ft/\partial\param_i}{\Jf}
    \;+\;
    \frac{1}{\Jf}\frac{\partial\Jf}{\partial\param_i}
  \right],
\end{equation}
where $s(z) = \partial_z\log\pz(z)$ is the \emph{score function} of the
latent density.  Analytical expressions for $s(z)$ are available for all
four latent families.

\subsection{Parameter Optimisers}

Four first-order optimisers are available.  In all cases $g$ denotes the
current (mini-batch) gradient, $\eta$ the learning rate, and $t$ the epoch
counter.

\paragraph{Stochastic Gradient Descent (SGD).}
The simplest update rule: move the parameters a fixed step $\eta$ in the
negative gradient direction,
\begin{equation}
  \param \;\leftarrow\; \param - \eta\,g.
\end{equation}
Each update uses only the current gradient with no memory of past
iterates.  Convergence can be slow and oscillatory near valleys where the
loss surface curves differently in different directions.

\paragraph{SGD with momentum.}
Accumulates a velocity vector $m$ as an exponential moving average of
past gradients,
\begin{equation}
  m \;\leftarrow\; \beta_1\,m + g,\qquad
  \param \;\leftarrow\; \param - \eta\,m.
\end{equation}
The momentum term ($\beta_1 = 0.9$) damps oscillations and accelerates
convergence along directions of consistent gradient sign, analogously to a
ball rolling with inertia.

\paragraph{RMSprop.}
Maintains a per-parameter running average of squared gradients,
\begin{equation}
  v \;\leftarrow\; \beta_2\,v + (1-\beta_2)\,g^2,\qquad
  \param \;\leftarrow\; \param
         - \frac{\eta}{\sqrt{v+\varepsilon}}\,g.
\end{equation}
Dividing by $\sqrt{v}$ normalises each parameter's effective step size by
the recent gradient magnitude, so parameters with large, noisy gradients
take smaller steps while those with small gradients take larger ones.
This adaptive scaling is particularly helpful when parameters operate on
very different scales.  The app uses $\beta_2=0.9$, $\varepsilon=10^{-8}$.

\paragraph{Adam (default).}
Combines momentum and per-parameter adaptive scaling with
bias-correction for the first few steps:
\begin{align}
  m &\;\leftarrow\; \beta_1\,m + (1-\beta_1)\,g,&
  \hat m &= \frac{m}{1-\beta_1^t},\\
  v &\;\leftarrow\; \beta_2\,v + (1-\beta_2)\,g^2,&
  \hat v &= \frac{v}{1-\beta_2^t},\\
  \param &\;\leftarrow\; \param
         - \frac{\eta\,\hat m}{\sqrt{\hat v}+\varepsilon}.
\end{align}
The bias-corrected estimates $\hat m$ and $\hat v$ prevent the
artificially small steps that would otherwise occur in the first few
epochs when $m$ and $v$ are initialised at zero.  Adam inherits both the
directional smoothing of momentum and the per-parameter scale adaptation
of RMSprop, making it robust to a wide range of loss landscape shapes and
the recommended default.  The app uses $\beta_1=0.9$, $\beta_2=0.999$,
$\varepsilon=10^{-8}$.

%% ─────────────────────────────────────────────────────────────────────────────
\section{Numerical Methods}

\subsection{Bisection Inversion}
\label{sec:bisection}

The NLL gradient requires evaluating $z = \ft^{-1}(x)$ for every data
point at every epoch.  For the SLP family this is done by bisection on the
bracketing interval $[-20, 20]$: since $\ft$ is strictly monotone, the
midpoint rule halves the interval at each step.  Starting from a bracket
of width 40, after $n$ steps the residual is at most $40/2^n$; 30 steps
reduce it below $4\times10^{-8}$, well within the target tolerance of
$10^{-7}$.

To avoid paying the bisection cost twice per epoch, the loss and gradient
are computed in a single pass: the inverse $z = \ft^{-1}(x)$ is found
once, the sigmoid activations and Jacobian are accumulated simultaneously,
and both the loss and its gradient are evaluated from those shared
intermediate quantities.

\subsection{Example-Data Generation}

Exact samples from the Boltzmann target $p^*$ are generated by
quantile-transform sampling:
\begin{equation}
  x_i^* \;=\; F_{p^*}^{-1}\!\left(F_{\pz}(z_i)\right),
\end{equation}
where $z_i \sim \pz$.  Both CDFs are computed numerically on fine grids
and the inversions are performed by linear interpolation.  This produces
samples that are exact (up to grid resolution) without requiring Markov
chain Monte Carlo.

%% ─────────────────────────────────────────────────────────────────────────────
\section{Application Interface}

\subsection{Layout}

The browser window is divided into two regions:
\begin{itemize}
  \item \textbf{Left — control panel.}  A fixed-width panel with a header
        image and three tabs (\emph{Densities}, \emph{Map}, \emph{Training}),
        followed by an auto-rescale checkbox and Reset / Exit buttons.  The
        tabs are shown as a segmented control with icons; the \emph{Exit}
        button closes the current browser session (it does not stop the
        server for other users).
  \item \textbf{Right — figure.}  A Matplotlib figure split into two
        columns: the main four-panel display (left) and three histogram /
        loss panels (right).
\end{itemize}

\subsection{Main Four-Panel Display}

The panels share their horizontal axis ($z$) and are arranged vertically:

\begin{center}
\renewcommand{\arraystretch}{1.3}
\begin{tabular}{lll}
\toprule
Panel & Horizontal axis & Vertical axis \\
\midrule
Top strip  & $z$ & $\pz(z)$ (latent density, blue) \\
Main panel & $z$ & $x = \ft(z)$ (transformation curve, red/orange) \\
Bottom strip & $z$ & $|\Jf(z)|^{-1}$ (inverse Jacobian, red/orange) \\
Right strip  & $\px(x)$ (density, green) & $x$ (shared with main) \\
\bottomrule
\end{tabular}
\end{center}

Both axes of the main panel carry equal padding of 12\,\% of the
respective data range beyond the displayed support, so the L-shaped
mapping lines (see below) have symmetric margins on both sides.

The main panel also shows the target density $p^*(x)$ (amber, when
\emph{Show target} is ticked) and the exact transformation $x^*(z)$
(dashed amber, when \emph{Show exact transformation} is ticked).
Regions where $\Jf < 0$ are drawn in orange to flag non-invertibility.
Enabling \emph{Show target} extends the $x$-range only if the target is
\emph{wider} than the transformed distribution; a narrower target leaves
the range unchanged (it never causes a spurious rescale).

\paragraph{Mapping lines.}
When \emph{Show mapping lines} is enabled (Densities tab), $N$ evenly
spaced values $z_1,\ldots,z_N$ are chosen across the displayed $z$ range
and each is connected to its image $x_i = \ft(z_i)$ by an L-shaped pair
of grey segments:
\begin{itemize}
  \item a \emph{vertical leg} descending from the top edge of the main
        panel down to the curve at $(z_i, x_i)$, and
  \item a \emph{horizontal leg} running from $(z_i, x_i)$ to the right
        edge of the main panel.
\end{itemize}
Lines end exactly at the panel borders; no markers are placed on the
density curves or on the transformation curve itself.  While mapping
lines are shown, the horizontal and vertical zero-reference lines in the
main panel are suppressed to reduce visual clutter.  The number $N$ can
be set between 1 and 100; values above 100 are automatically clamped.

\subsection{Densities Tab}

Controls the latent distribution (type and parameters) and the target
Boltzmann distribution (temperature $\kT$, potential coefficients
$u_1$--$u_4$).

The $\mu$ and $\sigma$ sliders behave as follows:
\begin{itemize}
  \item For \textbf{Gaussian, Uniform, Laplace, Cauchy}: $\mu$ is the
        mean (range $[-2,2]$, default $0$) and $\sigma$ is the scale
        parameter (range $[0.1,2]$, default $1$).
  \item For \textbf{Bimodal}: $\mu$ is the \emph{peak position} — the
        two peaks sit at $-\mu$ and $+\mu$ — so only positive values are
        meaningful.  The slider range is therefore restricted to $[0,2]$
        with default $\mu=1$.  $\sigma$ is the width of each individual
        peak (default $0.5$).  These defaults are applied automatically
        when Bimodal is selected from the drop-down.
\end{itemize}

\textbf{Show mapping lines} (checkbox + \emph{N\,=} entry) — overlays
$N \in [1, 100]$ L-shaped mapping lines in the main panel as described in
Section~5.2.  Entering a value larger than 100 snaps the field back to
100 on the next redraw.

\subsection{Map Tab}

Selects the transformation family (Polynomial, SLP, or RQS) and exposes
its parameters via sliders.  \emph{Defaults} resets to the near-identity
configuration; \emph{Randomise} draws random parameter values.

For the RQS, the number of bins $K\in\{2,3,4\}$ is chosen via radio
buttons.  The sliders expose the unconstrained logits for the bin widths,
bin heights, and knot derivatives; the actual constrained values are
computed via softmax and exp internally (see Section~\ref{sec:rqs}).
All logits set to zero produce the identity transformation within the
support $[-B,B]$.

\subsection{Training Tab}

\begin{itemize}
  \item \textbf{Sample!} — draws $N$ samples from $\pz$ and displays
        their histogram and the transformed histogram.
  \item \textbf{Generate data} — produces $N$ exact samples from $p^*$
        via quantile-transform sampling, for use with example-based
        training.
  \item \textbf{Mode} — selects Energy-based or Example-based training.
  \item \textbf{Train!} — runs the optimiser in a background thread for
        the specified number of epochs.  A timer animates the transformation
        in real time at a fixed 30\,ms refresh rate.  When training finishes
        the figure shows the \emph{best-loss} parameters encountered during
        the run (the lowest point of the loss curve), rather than the last
        iterate — so the final density is never worse than what was displayed
        while training, even when the loss oscillates near convergence.
  \item \textbf{Epochs / lr / N\,batch} — total epochs, learning rate,
        and mini-batch size (Energy-based) or full dataset size
        (Example-based).
  \item \textbf{Every $n$ ep} — stride between display updates; setting
        this to a larger value speeds up training at the cost of less
        frequent animation.
  \item \textbf{Post-training display.}  Sliders have finite ranges and
        silently clamp any trained parameter that falls outside them.  To
        prevent the displayed density from jumping after training ends, the
        app retains the full-precision parameter vector internally and uses
        it for all rendering as long as no slider has been manually adjusted
        since the last training run.  Moving any slider or resetting the
        parameters switches the display back to the slider values.
\end{itemize}

\subsection{Right-Column Panels}

\begin{itemize}
  \item \textbf{Latent $z$ histogram} — distribution of the current
        sample / training batch in $z$-space, overlaid with the
        theoretical $\pz$.  Its horizontal axis spans the same $z$-range as
        the top/main/inverse-Jacobian panels.
  \item \textbf{Transformed $x$ histogram} — same samples pushed through
        $\ft$, overlaid with the theoretical $\px$ (and the target $p^*$
        if \emph{Show target} is on).  Its horizontal axis spans the same
        $x$-range as the main panel's vertical axis.  When \emph{Show
        importance weights} is also ticked, the panel shows the
        importance-weight overlay described in Section~\ref{sec:iw}.
  \item \textbf{Training loss} — the total loss and its two components
        (energy/latent and entropy/Jacobian terms), drawn in three clearly
        distinct colours.  A fixed subsampling stride (set once at the start
        of each run) ensures previously plotted points never shift position
        as new epochs are added.
\end{itemize}

The density curves overlaid on both histograms (and on the side panels) are
drawn without clipping, so their line thickness stays constant even where a
curve meets the axis at density zero.

\subsection{Importance-Weight Overlay}
\label{sec:iw}

When both \emph{Show target} and \emph{Show importance weights} (Training
tab) are ticked and samples are present, the transformed-$x$ histogram panel
gains a secondary $y$-axis (right, brown) that shows the
\emph{unnormalised importance weights}
\begin{equation}
  w(x) \;=\; \frac{p^*(x)}{\px(x)}
           \;=\; \frac{p^*(x)\,|\Jf(z)|}{p_z(z)},
  \qquad z = \ft^{-1}(x).
  \label{eq:iw}
\end{equation}
A dashed brown curve traces $w(x)$ over the same $x$-range as the
density plot; a dotted reference line marks $w = 1$ (perfect agreement
between model and target).  Where $w > 1$ the model underestimates the
target; where $w < 1$ it overestimates it.

The \textbf{effective sample size}
\begin{equation}
  N_{\mathrm{eff}} \;=\;
  \frac{\left(\displaystyle\sum_{i=1}^N w_i\right)^{\!2}}
       {\displaystyle\sum_{i=1}^N w_i^2}
  \;\in\; (0, N],
  \label{eq:neff}
\end{equation}
where $w_i = p^*(x_i)/\px(x_i)$ at the actual sample points, is shown
alongside the \emph{mean importance weight} $\bar w = N^{-1}\sum_i w_i$
in the panel annotation.

\paragraph{Why $N_{\mathrm{eff}}$ alone is insufficient.}
$N_{\mathrm{eff}}$ measures \emph{weight uniformity}: it equals $N$ when
all $w_i$ are identical, regardless of their magnitude.
In particular, \emph{mode collapse} — where $\px$ concentrates on one
mode of a multimodal $p^*$ — produces weights that are all small but
equal, so $N_{\mathrm{eff}} \approx N$ even though the model is far from
the target.

\paragraph{Mean weight $\bar w$ as an absolute quality indicator.}
By the change-of-variables identity
$\mathbb{E}_{p_x}[w(x)] = \int p^*(x)\,dx = 1$,
the sample mean $\bar w$ should be close to $1$ for a well-fitted model.
With mode collapse, samples land where $\px \gg p^*$, giving
$w_i = p^*/\px \ll 1$ uniformly, so $\bar w \ll 1$.

The two metrics are complementary:
\begin{itemize}
  \item $\bar w \approx 1$: model covers the target correctly (absolute).
  \item $N_{\mathrm{eff}}/N \approx 100\,\%$: no single sample dominates
        (relative).  A low value flags over-dispersion.
\end{itemize}
During training, both should approach their ideal values ($\bar w\to 1$,
$N_{\mathrm{eff}}/N \to 100\,\%$) as the flow converges.

\subsection{Auto-Rescale Checkbox}

When ticked (default), the $z$- and $x$-axis ranges of all panels are
recomputed on every redraw to fit the current data.  When unticked, those
ranges are frozen at their current values, allowing close inspection of a
particular region during training or parameter sweeps.  The density axes
(heights of $\pz$, $|\Jf|^{-1}$, $\px$) continue to rescale freely in
both modes.

The axis ranges covered by auto-rescale are: the shared $z$-axis of the
top/main/bottom strips; the shared $x$-axis of the main and right strips;
the $z$-axis of the latent histogram; and the $x$-axis of the transformed
histogram.  The transformed-histogram $x$-axis is computed explicitly from
the union of all plotted data ranges (samples, target curve, example data)
with 4\,\% padding, ensuring it re-fits correctly when auto-rescale is
re-enabled after having been frozen.

%% ─────────────────────────────────────────────────────────────────────────────
\begin{thebibliography}{9}
\bibitem{rezende2015}
D.~J. Rezende and S.~Mohamed,
\newblock Variational inference with normalizing flows,
\newblock \emph{Proceedings of the 32nd International Conference on Machine
Learning (ICML)}, PMLR 37:1530--1538, 2015.

\bibitem{papamakarios2021}
G.~Papamakarios, E.~Nalisnick, D.~J. Rezende, S.~Mohamed, and
B.~Lakshminarayanan,
\newblock Normalizing flows for probabilistic modeling and inference,
\newblock \emph{Journal of Machine Learning Research} \textbf{22}(57):1--64,
2021.

\bibitem{noe2019boltzmann}
F.~No\'e, S.~Olsson, J.~K\"ohler, and H.~Wu,
\newblock Boltzmann generators: Sampling equilibrium states of many-body
systems with deep learning,
\newblock \emph{Science} \textbf{365}, eaaw1147, 2019.

\bibitem{coretti2024kimreview}
A.~Coretti, S.~Falkner, J.~Weinreich, C.~Dellago, and O.~A. von~Lilienfeld,
\newblock Boltzmann Generators and the new frontier of computational sampling
in many-body systems,
\newblock \emph{KIM Review} \textbf{2}, Article 03, 2024.

\bibitem{wirnsberger2020}
P.~Wirnsberger, A.~J. Ballard, G.~Papamakarios, S.~Abercrombie,
S.~Racani\`ere, A.~Pritzel, D.~Jimenez~Rezende, and C.~Blundell,
\newblock Targeted free energy estimation via learned mappings,
\newblock \emph{The Journal of Chemical Physics} \textbf{153}, 144112, 2020.

\bibitem{falkner2023}
S.~Falkner, A.~Coretti, S.~Romano, P.~L. Geissler, and C.~Dellago,
\newblock Conditioning Boltzmann Generators for rare event sampling,
\newblock \emph{Machine Learning: Science and Technology} \textbf{4}, 035050,
2023.

\bibitem{schebek2024}
M.~Schebek, M.~Invernizzi, F.~No\'e, and J.~Rogal,
\newblock Efficient mapping of phase diagrams with conditional Boltzmann
Generators,
\newblock \emph{Machine Learning: Science and Technology} \textbf{5}, 045045,
2024.

\bibitem{schebek2025}
M.~Schebek, F.~No\'e, and J.~Rogal,
\newblock Scalable Boltzmann Generators for equilibrium sampling of
large-scale materials,
\newblock \emph{arXiv:2509.25486}, 2025.

\bibitem{coretti2025jcp}
A.~Coretti, S.~Falkner, P.~L. Geissler, and C.~Dellago,
\newblock Learning mappings between equilibrium states of liquid systems using
normalizing flows,
\newblock \emph{The Journal of Chemical Physics} \textbf{162}, 184102, 2025.

\bibitem{durkan2019neural}
C.~Durkan, A.~Bekasov, I.~Murray, and G.~Papamakarios,
\newblock Neural spline flows,
\newblock \emph{Advances in Neural Information Processing Systems}, 32, 2019.
\end{thebibliography}

\end{document}
